% Revisione Ready
% Grammatica Ready
% Citazioni Ready

\section{Grafi e Ipergrafi}\label{sec:grafi_e_ipergrafi}

Quando all'interno del seguente lavoro si trattano i \textit{grafi}, ci si riferisce ad una struttura matematica composta da un insieme di nodi (o vertici) e da un insieme di archi che connettono coppie di nodi.
Formalmente, un grafo è una coppia $G = (V, E)$, dove $V$ è l'insieme di nodi e $E$ è l'insieme di coppie di nodi, definita formalmente come:
\[
    E \subseteq \{\{u,v\} \mid u,v \in V,\; u \neq v\}.
\]
Inoltre, il grado di un nodo $v$, indicato con $\deg(v)$, rappresenta il numero di archi incidenti su di esso.
I grafi si possono poi classificare in base a diverse caratteristiche, come ad esempio l'orientamento degli archi o la presenza di un fattore di peso che può essere associato agli archi o ai nodi.
I grafi costituiscono uno strumento fondamentale per la modellazione di sistemi complessi e trovano applicazione nell'analisi delle reti sociali, delle infrastrutture di trasporto, delle reti di comunicazione e dei sistemi biologici~\cite{newman2018}.\newline
Al contempo, però, tali strutture dati possono risultare estremamente limitanti quando si tratta di rappresentare, senza alcuna perdita di informazione, interazioni che coinvolgono molteplici entità contemporaneamente.\newline
Gli \textit{ipergrafi} sono invece la generalizzazione dei grafi; tali strutture dati complesse sono composte da un insieme di nodi e da un insieme di iperarchi, all'interno dei quale ogni iperarco può connettere un numero arbitrario di nodi.
Formalmente, un ipergrafo è una coppia $H = (V, E)$, dove $V$ è l'insieme di nodi e $E$ è l'insieme di sottoinsiemi di $V$, definito formalmente come:
\[
    E \subseteq \{ e \subseteq V \mid e \neq \emptyset \}.
\]
La definizione di grado non cambia particolarmente per gli ipergrafi; tuttavia, introducendo la cardinalità di un iperarco, indicata con $|e|$, si rappresenta il numero di nodi coinvolti nell'interazione descritta dall'iperarco~\cite{bretto2013}.
Come per i grafi, anche gli ipergrafi possono essere classificati in base a diverse caratteristiche, come ad esempio l'orientamento degli iperarchi o la presenza di un fattore di peso che può essere associato agli iperarchi o ai nodi.
Tuttavia, all'interno di questo lavoro verranno considerati solamente ipergrafi non orientati e senza peso.
Gli ipergrafi, a differenza dei grafi, consentono di rappresentare senza perdita di informazione interazioni che coinvolgono molteplici entità contemporaneamente.
Tale generalizzazione permette di rappresentare sistemi caratterizzati da interazioni collettive~\cite{battiston2020, antelmi2023}.